% ID: SEL-D0058-OA-E01
% ESTADO: validada. S058.
\textbf{a)} En el primer cuadrante, la recta $x+y=2$ recorta el triángulo junto al origen. Sobre $y=0$ queda $2\le x\le5$, y sobre $x=0$ queda $2\le y\le5$. Las rectas $x+3y=15$ y $3x-y=15$ se cruzan en $(6,3)$, que cumple todas las restricciones.
\begin{center}
\setlength{\unitlength}{0.9mm}
\begin{picture}(92,65)(-12,-8)
\qbezier(0,0)(34,0)(68,0)
\qbezier(0,0)(0,25)(0,50)
\put(69,0){$x$}\put(2,50){$y$}
\qbezier(20,0)(35,0)(50,0)
\qbezier(50,0)(55,12)(60,24)
\qbezier(60,24)(30,32)(0,40)
\qbezier(0,40)(0,28)(0,16)
\qbezier(0,16)(10,8)(20,0)
\put(17,-6){\scriptsize $(2,0)$}\put(47,-6){\scriptsize $(5,0)$}
\put(62,24){\scriptsize $(6,3)$}\put(2,42){\scriptsize $(0,5)$}
\put(-12,17){\scriptsize $(0,2)$}
\end{picture}
\end{center}
\textbf{b)} En los cinco vértices, los valores son
\[
\begin{array}{c|rrrrr}
(x,y)&(2,0)&(5,0)&(6,3)&(0,5)&(0,2)\\\hline
F&6&15&21&5&2
\end{array}
\]
Por linealidad, el máximo es $21$ en $(6,3)$ y el mínimo $2$ en $(0,2)$. Cada uno se alcanza sólo en su vértice, pues los demás vértices dan valores distintos.

\textbf{c)} Como $30>21$, no hay ningún punto factible con ese valor.
